\section{Topología euclídea}

\begin{definition}[Norma euclídea]
  La \emph{norma euclídea} en $\R^n$ es la función $\norm{\cdot}\colon \R^n \to \R$ dada por
  \[
    \norm{x} = \sqrt{\sum_{i=1}^{n} x_i^2}, \qquad x = (x_1, \dots, x_n).
  \]
\end{definition}

Se cumplen las siguientes propiedades, para todo $x, y \in \R^n$ y todo $\lambda \in \R$:
\begin{enumerate}[label=(\arabic*)]
  \item $\norm{x} \geq 0$, y $\norm{x} = 0 \iff x = 0$.
  \item $\norm{\lambda x} = |\lambda| \, \norm{x}$.
  \item $\norm{x + y} \leq \norm{x} + \norm{y}$ (desigualdad triangular).
\end{enumerate}
Las dos primeras son inmediatas; la tercera se deduce de la siguiente desigualdad.

\begin{prop}[Desigualdad de Cauchy--Schwarz]
  Dados $x, y \in \R^n$,
  \[
    \left| \sum_{i=1}^{n} x_i y_i \right| \leq \norm{x} \, \norm{y}.
  \]
\end{prop}

\begin{proof}
  Se verifica la identidad de Lagrange
  \[
    \left( \sum_{k=1}^{n} x_k^2 \right) \left( \sum_{k=1}^{n} y_k^2 \right) - {\left( \sum_{k=1}^{n} x_k y_k \right)}^2 = \sum_{1 \leq k < j \leq n} {(x_k y_j - x_j y_k)}^2.
  \]
  En efecto, el lado izquierdo es
  \[
    \sum_{k, j} x_k^2 y_j^2 - \sum_{k, j} x_k y_k x_j y_j = \frac{1}{2} \sum_{k, j} \bigl( x_k^2 y_j^2 + x_j^2 y_k^2 - 2 x_k y_k x_j y_j \bigr) = \frac{1}{2} \sum_{k, j} {(x_k y_j - x_j y_k)}^2,
  \]
  y los términos con $k = j$ se anulan, mientras que los pares $(k, j)$ y $(j, k)$ aportan lo mismo. Como el lado derecho es una suma de cuadrados, es no negativo, de modo que ${(\sum x_k y_k)}^2 \leq \norm{x}^2 \norm{y}^2$.
\end{proof}

\begin{proof}[Demostración de la desigualdad triangular]
  Por Cauchy--Schwarz,
  \[
    \norm{x + y}^2 = \norm{x}^2 + 2 \sum_i x_i y_i + \norm{y}^2 \leq \norm{x}^2 + 2 \norm{x} \norm{y} + \norm{y}^2 = {\bigl( \norm{x} + \norm{y} \bigr)}^2. \qedhere
  \]
\end{proof}

\begin{definition}[Distancia euclídea]
  La \emph{distancia euclídea} es la función $d\colon \R^n \times \R^n \to \R$, $d(x, y) = \norm{x - y}$.
\end{definition}

\begin{prop}
  Para todos $x, y, z \in \R^n$:
  \begin{enumerate}[label=(\arabic*)]
    \item $d(x, y) \geq 0$, y $d(x, y) = 0 \iff x = y$.
    \item $d(x, y) = d(y, x)$.
    \item $d(x, z) \leq d(x, y) + d(y, z)$.
  \end{enumerate}
\end{prop}

\begin{proof}
  (1) Es la primera propiedad de la norma. (2) $d(x, y) = \norm{x - y} = \norm{(-1)(y - x)} = \norm{y - x} = d(y, x)$. (3) $d(x, z) = \norm{(x - y) + (y - z)} \leq \norm{x - y} + \norm{y - z} = d(x, y) + d(y, z)$.
\end{proof}

\section{Conjuntos abiertos}

\begin{definition}[Bola abierta]
  Si $x \in \R^n$ y $r > 0$, la \emph{bola abierta} de centro $x$ y radio $r$ es el conjunto
  \[
    B_r(x) = B(x, r) = \{ y \in \R^n : d(x, y) < r \}.
  \]
\end{definition}

\begin{definition}[Conjunto abierto]
  Un subconjunto $U \subset \R^n$ es \emph{abierto} si para todo $x \in U$ existe $r > 0$ tal que $B_r(x) \subset U$.
\end{definition}

\begin{lemma}
  La bola abierta $B_r(x)$ es un conjunto abierto.
\end{lemma}

\begin{proof}
  Sea $y \in B_r(x)$ y $s = r - d(x, y) > 0$. Si $z \in B_s(y)$, entonces
  \[
    d(x, z) \leq d(x, y) + d(y, z) < d(x, y) + s = r,
  \]
  es decir, $z \in B_r(x)$. Luego $B_s(y) \subset B_r(x)$.
\end{proof}

\begin{prop}
  Un subconjunto $U \subsetneq \R^n$ es abierto si y sólo si para todo $x \in U$ se tiene $\inf \{ d(x, y) : y \in \R^n \setminus U \} > 0$. El ínfimo existe pues el conjunto es no vacío y acotado inferiormente por $0$.
\end{prop}

\begin{proof}
  ($\Rightarrow$) Sea $x \in U$ y $r > 0$ con $B_r(x) \subset U$. Si $y \in \R^n \setminus U$, entonces $y \notin B_r(x)$, es decir, $d(x, y) \geq r$. Luego $r$ es cota inferior del conjunto y el ínfimo es mayor o igual que $r > 0$.

  ($\Leftarrow$) Sea $x \in U$ y $s = \inf \{ d(x, y) : y \in \R^n \setminus U \} > 0$. Si $z \in B_s(x)$ y fuera $z \notin U$, sería $d(x, z) \geq s$, lo que contradice $d(x, z) < s$. Luego $B_s(x) \subset U$.
\end{proof}

\begin{prop}
  Valen las siguientes propiedades:
  \begin{enumerate}[label=(\arabic*)]
    \item $\varnothing$ y $\R^n$ son abiertos.
    \item La intersección de una familia finita y no vacía de abiertos de $\R^n$ es abierta.
    \item La unión de una familia cualquiera de abiertos de $\R^n$ es abierta.
  \end{enumerate}
\end{prop}

\begin{proof}
  (1) Es inmediato de la definición.

  (2) Sean $U_1, \dots, U_k$ abiertos y $x \in U = \bigcap_{i=1}^{k} U_i$. Para cada $i$ existe $r_i > 0$ con $B_{r_i}(x) \subset U_i$. Con $r = \min\{r_1, \dots, r_k\} > 0$ se tiene $B_r(x) \subset U_i$ para todo $i$, es decir, $B_r(x) \subset U$.

  (3) Sea $\mathscr{U}$ una familia de abiertos y $U$ su unión. Dado $x \in U$, existe $V \in \mathscr{U}$ con $x \in V$. Como $V$ es abierto, existe $r > 0$ con $B_r(x) \subset V \subset U$.
\end{proof}

La hipótesis de que la familia sea finita en (2) es necesaria: si $U_n = \left( 1 - \frac{1}{n}, 1 + \frac{1}{n} \right)$ para $n \in \N$, entonces $\bigcap_{n \geq 1} U_n = \{1\}$, que no es abierto.

\begin{definition}[Topología euclídea]
  La \emph{topología euclídea} de $\R^n$ es el conjunto de todos los abiertos de $\R^n$.
\end{definition}

\begin{corollary}
  Un subconjunto de $\R^n$ es abierto si y sólo si es unión de una familia de bolas abiertas.
\end{corollary}

\begin{proof}
  ($\Rightarrow$) Sea $U$ abierto y sea $\mathscr{B}$ la familia de bolas abiertas contenidas en $U$, con unión $\tilde{U}$. Claramente $\tilde{U} \subset U$. Recíprocamente, dado $x \in U$, existe $r > 0$ con $B_r(x) \subset U$, es decir, $B_r(x) \in \mathscr{B}$, y entonces $x \in B_r(x) \subset \tilde{U}$. Luego $U = \tilde{U}$.

  ($\Leftarrow$) Las bolas abiertas son abiertas, y toda unión de abiertos es abierta.
\end{proof}

\section{Interior de un conjunto}

\begin{definition}[Interior]
  Sea $X \subset \R^n$. Un punto $x \in \R^n$ es \emph{interior} a $X$ si existe $r > 0$ con $B_r(x) \subset X$. El \emph{interior} de $X$, que notamos $X^{\circ}$, es el conjunto de los puntos interiores de $X$.
\end{definition}

\begin{note}
  Siempre vale $X^{\circ} \subset X$, y $X$ es abierto si y sólo si $X^{\circ} = X$.
\end{note}

\begin{prop}
  Sea $X \subset \R^n$. Entonces $X^{\circ}$ es la unión de todos los subconjuntos abiertos de $X$; en particular es abierto, y es el mayor abierto contenido en $X$.
\end{prop}

\begin{proof}
  Sea $U$ la unión de todos los abiertos contenidos en $X$. Veamos que $U = X^{\circ}$.

  ($X^{\circ} \subset U$) Si $x \in X^{\circ}$, existe $r > 0$ con $B_r(x) \subset X$; como $B_r(x)$ es abierto, $B_r(x) \subset U$, y entonces $x \in U$.

  ($U \subset X^{\circ}$) Si $x \in U$, existe un abierto $V \subset X$ con $x \in V$, y existe $r > 0$ con $B_r(x) \subset V \subset X$. Luego $x \in X^{\circ}$.

  Como $U$ es unión de abiertos, es abierto, y contiene a todo abierto contenido en $X$.
\end{proof}

\begin{corollary}
  Si $X \subset Y$, entonces $X^{\circ} \subset Y^{\circ}$.
\end{corollary}

\begin{proof}
  $X^{\circ}$ es un abierto contenido en $X \subset Y$, y $Y^{\circ}$ es el mayor abierto contenido en $Y$.
\end{proof}

\subsection*{Propiedades del interior}

\begin{prop}
  ${(X \cap Y)}^{\circ} = X^{\circ} \cap Y^{\circ}$.
\end{prop}

\begin{proof}
  Como $X^{\circ} \cap Y^{\circ}$ es abierto (intersección finita de abiertos) y está contenido en $X \cap Y$, se tiene $X^{\circ} \cap Y^{\circ} \subset {(X \cap Y)}^{\circ}$. Por otro lado, $X \cap Y \subset X$ implica ${(X \cap Y)}^{\circ} \subset X^{\circ}$, y análogamente ${(X \cap Y)}^{\circ} \subset Y^{\circ}$. Luego ${(X \cap Y)}^{\circ} \subset X^{\circ} \cap Y^{\circ}$.
\end{proof}

\begin{prop}
  $X^{\circ} \cup Y^{\circ} \subset {(X \cup Y)}^{\circ}$.
\end{prop}

\begin{proof}
  $X^{\circ} \cup Y^{\circ}$ es abierto y está contenido en $X \cup Y$, de modo que está contenido en el mayor abierto contenido en $X \cup Y$.
\end{proof}

\begin{eg}
  La inclusión anterior puede ser estricta. Sean $X = (-1, 1)$ e $Y = [1, 2]$. Entonces $X^{\circ} \cup Y^{\circ} = (-1, 1) \cup (1, 2)$ no contiene a $1$, pero $X \cup Y = (-1, 2]$ tiene interior $(-1, 2) \ni 1$.
\end{eg}

\section{Entornos}

\begin{definition}[Entorno]
  Sea $x \in \R^n$. Un \emph{entorno} de $x$ es un subconjunto $N \subset \R^n$ tal que $x \in N^{\circ}$. Un \emph{entorno abierto} de $x$ es un entorno que es abierto, es decir, un abierto que contiene a $x$.
\end{definition}

\begin{prop}
  Un subconjunto de $\R^n$ es abierto si y sólo si es entorno de cada uno de sus puntos.
\end{prop}

\begin{proof}
  Si $U$ es abierto, $U = U^{\circ}$ y entonces $x \in U^{\circ}$ para todo $x \in U$. Recíprocamente, si $U$ es entorno de cada uno de sus puntos, $U \subset U^{\circ} \subset U$, de modo que $U = U^{\circ}$ y $U$ es abierto.
\end{proof}

\begin{prop}
  El interior de $X \subset \R^n$ es el conjunto de puntos de los que $X$ es entorno.
\end{prop}

\begin{proof}
  Es la definición: $X$ es entorno de $x$ si y sólo si $x \in X^{\circ}$.
\end{proof}

\section{Conjuntos cerrados}

\begin{definition}[Conjunto cerrado]
  Un subconjunto $F \subset \R^n$ es \emph{cerrado} si su complemento $\R^n \setminus F$ es abierto.
\end{definition}

\begin{prop}
  Valen las siguientes propiedades:
  \begin{enumerate}[label=(\arabic*)]
    \item $\varnothing$ y $\R^n$ son cerrados.
    \item La unión de una familia finita de cerrados es cerrada.
    \item La intersección de una familia arbitraria de cerrados es cerrada.
  \end{enumerate}
\end{prop}

\begin{proof}
  (1) $\R^n \setminus \R^n = \varnothing$ y $\R^n \setminus \varnothing = \R^n$ son abiertos.

  (2) Si $F_1, \dots, F_k$ son cerrados, $\R^n \setminus (F_1 \cup \cdots \cup F_k) = (\R^n \setminus F_1) \cap \cdots \cap (\R^n \setminus F_k)$ es una intersección finita de abiertos.

  (3) Si $\mathscr{F}$ es una familia de cerrados, $\R^n \setminus \bigcap_{F \in \mathscr{F}} F = \bigcup_{F \in \mathscr{F}} (\R^n \setminus F)$ es unión de abiertos.
\end{proof}

\begin{definition}[Bola cerrada]
  Si $x \in \R^n$ y $r > 0$, la \emph{bola cerrada} de centro $x$ y radio $r$ es
  \[
    \cl{B}_r(x) = \{ y \in \R^n : d(x, y) \leq r \}.
  \]
\end{definition}

\begin{prop}
  La bola cerrada $\cl{B}_r(x)$ es un conjunto cerrado.
\end{prop}

\begin{proof}
  Sea $y \in \R^n \setminus \cl{B}_r(x)$, es decir, $d(x, y) > r$, y sea $r_1 = d(x, y) - r > 0$. Si $z \in B_{r_1}(y)$, por la desigualdad triangular
  \[
    d(x, z) \geq d(x, y) - d(z, y) > d(x, y) - r_1 = r,
  \]
  de modo que $z \notin \cl{B}_r(x)$. Luego $B_{r_1}(y) \subset \R^n \setminus \cl{B}_r(x)$, y el complemento es abierto.
\end{proof}

\section{Clausura}

\begin{definition}[Clausura]
  La \emph{clausura} de $X \subset \R^n$ es la intersección $\cl{X}$ de todos los subconjuntos cerrados de $\R^n$ que contienen a $X$.
\end{definition}

\begin{prop}
  $\cl{X}$ es el menor subconjunto cerrado de $\R^n$ que contiene a $X$, en el sentido de que está contenido en todo cerrado que contiene a $X$.
\end{prop}

\begin{proof}
  $\cl{X}$ es cerrado por ser intersección de cerrados, contiene a $X$, y por definición está contenido en todo cerrado que contiene a $X$.
\end{proof}

\begin{corollary}
  Un subconjunto $X \subset \R^n$ es cerrado si y sólo si $\cl{X} = X$.
\end{corollary}

\begin{proof}
  Si $X$ es cerrado, es un cerrado que contiene a $X$, luego $\cl{X} \subset X$; como $X \subset \cl{X}$, vale la igualdad. Recíprocamente, si $X = \cl{X}$, entonces $X$ es cerrado pues $\cl{X}$ lo es.
\end{proof}

\begin{prop}
  Si $X \subset Y$, entonces $\cl{X} \subset \cl{Y}$.
\end{prop}

\begin{proof}
  $X \subset Y \subset \cl{Y}$, y $\cl{X}$ es el menor cerrado que contiene a $X$.
\end{proof}

\begin{prop}
  $\cl{X \cup Y} = \cl{X} \cup \cl{Y}$.
\end{prop}

\begin{proof}
  Como $X \subset X \cup Y$ e $Y \subset X \cup Y$, se tiene $\cl{X} \subset \cl{X \cup Y}$ y $\cl{Y} \subset \cl{X \cup Y}$, de donde $\cl{X} \cup \cl{Y} \subset \cl{X \cup Y}$. Por otro lado, $\cl{X} \cup \cl{Y}$ es cerrado (unión finita de cerrados) y contiene a $X \cup Y$, de modo que $\cl{X \cup Y} \subset \cl{X} \cup \cl{Y}$.
\end{proof}

\begin{prop}
  $\cl{X \cap Y} \subset \cl{X} \cap \cl{Y}$.
\end{prop}

\begin{proof}
  $\cl{X} \cap \cl{Y}$ es cerrado y contiene a $X \cap Y$, y $\cl{X \cap Y}$ es el menor cerrado que contiene a $X \cap Y$.
\end{proof}

\begin{prop}
  Sean $X \subset \R^n$ y $x \in \R^n$. Son equivalentes:
  \begin{enumerate}[label=(\arabic*)]
    \item $x \in \cl{X}$.
    \item Para todo $r > 0$, $B_r(x) \cap X \neq \varnothing$.
    \item Para todo entorno abierto $N$ de $x$, $N \cap X \neq \varnothing$.
    \item Para todo entorno $N$ de $x$, $N \cap X \neq \varnothing$.
  \end{enumerate}
\end{prop}

\begin{proof}
  (4) $\Rightarrow$ (3) $\Rightarrow$ (2) son inmediatos, pues los entornos abiertos son entornos y las bolas $B_r(x)$ son entornos abiertos de $x$.

  (2) $\Rightarrow$ (1) Por contrarrecíproco, supongamos $x \notin \cl{X}$. Entonces existe un cerrado $F$ con $X \subset F$ y $x \notin F$. Luego $x \in \R^n \setminus F$, que es abierto, y existe $r > 0$ con $B_r(x) \subset \R^n \setminus F$. Entonces $B_r(x) \cap X \subset (\R^n \setminus F) \cap F = \varnothing$.

  (1) $\Rightarrow$ (4) Por contrarrecíproco, supongamos que existe un entorno $N$ de $x$ con $N \cap X = \varnothing$. El conjunto $U = N^{\circ}$ es abierto, $x \in U$ y $U \subset N \subset \R^n \setminus X$. Entonces $X \subset \R^n \setminus U$, que es cerrado, de modo que $\cl{X} \subset \R^n \setminus U$ y, como $x \in U$, $x \notin \cl{X}$.
\end{proof}

\begin{eg}
  Sea $\Q^n$ el conjunto de puntos de $\R^n$ con todas sus coordenadas racionales. Entonces $\cl{\Q^n} = \R^n$; se dice que $\Q^n$ es \emph{denso} en $\R^n$.
\end{eg}

\begin{proof}
  Sea $x = (x_1, \dots, x_n) \in \R^n$ y $r > 0$. Para cada $i$, como $\Q$ es denso en $\R$, existe $q_i \in \Q$ con $|x_i - q_i| < \frac{r}{\sqrt{n}}$. Entonces $q = (q_1, \dots, q_n) \in \Q^n$ y
  \[
    d(x, q) = \sqrt{\sum_{i=1}^{n} {(x_i - q_i)}^2} < \sqrt{n \cdot \frac{r^2}{n}} = r,
  \]
  es decir, $q \in B_r(x) \cap \Q^n$. Por la equivalencia (1)~$\Leftrightarrow$~(2) de la proposición anterior, $x \in \cl{\Q^n}$.
\end{proof}
